\subsection{AI Agent}

Quá trình lựa chọn được thực hiện dựa trên việc so sánh các giải pháp phổ biến cho từng bài toán, với các tiêu chí chính bao gồm: mức độ phù hợp với yêu cầu hệ thống, khả năng mở rộng, độ phức tạp triển khai và chi phí vận hành.

Các công nghệ chính được sử dụng trong hệ thống bao gồm:

\begin{itemize}
\item PostgreSQL / Supabase (cơ sở dữ liệu quan hệ)
\item Neo4j / AuraDB (cơ sở dữ liệu đồ thị)
\item ChromaDB (cơ sở dữ liệu vector)
\item LangGraph (xây dựng AI Agent)
\item FastAPI (xây dựng API)
% \item Railway (triển khai hệ thống)
\end{itemize}

\subsubsection{Cơ sở dữ liệu quan hệ: PostgreSQL / Supabase}

Hệ thống cần một cơ sở dữ liệu để lưu trữ dữ liệu sản phẩm có cấu trúc và hỗ trợ các truy vấn lọc theo điều kiện như giá, thương hiệu và thông số kỹ thuật.

Các lựa chọn phổ biến được xem xét bao gồm:

\begin{tabular}{|p{2.5cm}|p{5.5cm}|p{5.5cm}|}
\hline
Công nghệ & Ưu điểm & Nhược điểm \\
\hline
MySQL & Phổ biến, dễ sử dụng & Hạn chế extension nâng cao \\
\hline
PostgreSQL & Hỗ trợ mạnh (index, extension, JSON) & Cấu hình phức tạp hơn \\
\hline
MongoDB & Linh hoạt schema & Không tối ưu cho truy vấn quan hệ \\
\hline
\end{tabular}

PostgreSQL được lựa chọn do khả năng hỗ trợ truy vấn phức tạp, hệ thống index mạnh và khả năng mở rộng thông qua extension. Những đặc điểm này phù hợp với yêu cầu xử lý dữ liệu sản phẩm có nhiều thuộc tính và cần lọc theo nhiều điều kiện.

Trong triển khai thực tế, PostgreSQL được sử dụng thông qua nền tảng Supabase. Việc sử dụng Supabase giúp tận dụng một dịch vụ database được quản lý sẵn, giảm đáng kể chi phí vận hành và cấu hình, đồng thời cung cấp các tính năng bổ trợ như quản lý dữ liệu và backup.

\subsubsection{Cơ sở dữ liệu đồ thị: Neo4j / AuraDB}

Đối với các bài toán liên quan đến quan hệ giữa các thực thể, cơ sở dữ liệu đồ thị là lựa chọn phù hợp hơn so với mô hình quan hệ truyền thống.

Các giải pháp được xem xét:

\begin{tabular}{|p{3cm}|p{5cm}|p{5cm}|}
\hline
Công nghệ & Ưu điểm & Nhược điểm \\
\hline
PostgreSQL (self-join) & Không cần thêm hệ thống mới & Query phức tạp, kém hiệu quả \\
\hline
Neo4j & Tối ưu cho graph, query trực quan (Cypher) & Cần học thêm ngôn ngữ \\
\hline
ArangoDB & Multi-model & Ít phổ biến hơn \\
\hline
\end{tabular}

Neo4j được lựa chọn do là hệ quản trị đồ thị phổ biến, hỗ trợ ngôn ngữ truy vấn Cypher trực quan và tối ưu cho các truy vấn quan hệ phức tạp.

Trong triển khai, hệ thống sử dụng Neo4j AuraDB — dịch vụ Neo4j được quản lý trên cloud. Điều này giúp loại bỏ nhu cầu tự vận hành database và đơn giản hóa quá trình triển khai.

\subsubsection{Cơ sở dữ liệu vector: ChromaDB}

Để hỗ trợ tìm kiếm ngữ nghĩa trong các câu hỏi tự nhiên, hệ thống cần một cơ sở dữ liệu vector có khả năng lưu trữ và truy vấn embedding.

Các lựa chọn phổ biến:

\begin{tabular}{|p{2.5cm}|p{5.5cm}|p{5.5cm}|}
\hline
Công nghệ & Ưu điểm & Nhược điểm \\
\hline
FAISS & Hiệu năng cao & Không có persistence mặc định \\
\hline
Pinecone & Managed, dễ dùng & Chi phí cao \\
\hline
ChromaDB & Nhẹ, dễ tích hợp, hỗ trợ local & Ít tính năng nâng cao \\
\hline
\end{tabular}

ChromaDB được lựa chọn do dễ tích hợp với Python và hỗ trợ lưu trữ cục bộ. Điều này cho phép hệ thống triển khai đơn giản mà không cần thêm một dịch vụ bên ngoài, đồng thời vẫn đáp ứng tốt nhu cầu tìm kiếm ngữ nghĩa với quy mô dữ liệu của đồ án.

\subsubsection{Công nghệ xây dựng AI Agent: LangGraph}

Hệ thống hội thoại yêu cầu khả năng điều phối giữa mô hình ngôn ngữ và các công cụ truy xuất dữ liệu theo nhiều bước. Do đó, cần một framework hỗ trợ orchestration và quản lý trạng thái.

Các lựa chọn:

\begin{tabular}{|p{3cm}|p{5cm}|p{5cm}|}
\hline
Công nghệ & Ưu điểm & Nhược điểm \\
\hline
Prompt chaining (thuần) & Đơn giản & Khó mở rộng, khó kiểm soát \\
\hline
LangChain Agent & Linh hoạt & Khó kiểm soát luồng \\
\hline
LangGraph & Kiểm soát tốt, hỗ trợ stateful & Phức tạp hơn \\
\hline
\end{tabular}

LangGraph được lựa chọn do cho phép xây dựng AI Agent dưới dạng đồ thị có trạng thái, giúp kiểm soát rõ ràng luồng xử lý và dễ dàng mở rộng.

\subsubsection{Mô hình ngôn ngữ lớn (LLM) và mô hình embedding}

AI Agent yêu cầu hai năng lực ngôn ngữ khác nhau: trích xuất thông tin có cấu trúc từ câu hỏi người dùng và sinh phản hồi tự nhiên kết hợp gọi công cụ. Do đặc thù của từng bài toán, hệ thống sử dụng hai mô hình ngôn ngữ riêng biệt được tích hợp qua API, không thực hiện huấn luyện hay fine-tuning.

Cần làm rõ rằng đây là kiến trúc RAG kết hợp Tool-calling Agent: các mô hình LLM được sử dụng hoàn toàn như một dịch vụ suy luận (inference service). Kiến thức về sản phẩm và chính sách không được đưa vào mô hình qua quá trình huấn luyện mà được truy xuất động từ cơ sở dữ liệu tại thời điểm xử lý yêu cầu. Điều này giúp hệ thống cập nhật dữ liệu mà không cần huấn luyện lại mô hình.

Node Parse yêu cầu mô hình có khả năng sinh ra đầu ra tuân theo schema cố định (structured output). Các lựa chọn được xem xét:

\begin{table}[H]
    \centering
    \small
    \renewcommand{\arraystretch}{1.3}
    \begin{tabular}{|p{3cm}|p{5cm}|p{5cm}|}
        \hline
        \textbf{Mô hình} & \textbf{Ưu điểm} & \textbf{Nhược điểm} \\
        \hline
        
        GPT-4o (OpenAI)
        & Structured output chính xác, chất lượng cao, hỗ trợ tiếng Việt tốt
        & Chi phí cao hơn so với các mô hình nhỏ hơn \\
        \hline
        
        GPT-4o mini (OpenAI)
        & Structured output ổn định, tốc độ nhanh hơn GPT-4o, chi phí thấp, hỗ trợ tiếng Việt tốt
        & Chất lượng suy luận phức tạp kém hơn GPT-4o \\
        \hline
        
        Gemini 2.5 Flash (Google)
        & Tốc độ nhanh, chi phí thấp
        & Structured output kém ổn định hơn với schema phức tạp trong thực nghiệm \\
        \hline
    \end{tabular}
\end{table}

GPT-4o mini được lựa chọn vì cung cấp structured output ổn định với schema phức tạp (nhiều trường lồng nhau), tốc độ phản hồi nhanh và chi phí thấp hơn GPT-4o. Node Parse được gọi ở mỗi lượt hội thoại nên tốc độ và chi phí là yếu tố quan trọng, trong khi chất lượng structured output vẫn đáp ứng đủ yêu cầu.

Node Reason yêu cầu mô hình có khả năng lập luận phức tạp, gọi công cụ (tool calling) linh hoạt và sinh phản hồi tự nhiên bằng tiếng Việt. Các lựa chọn được xem xét:

\begin{table}[H]
    \centering
    \small
    \renewcommand{\arraystretch}{1.3}
    \begin{tabular}{|p{3cm}|p{5cm}|p{5cm}|}
        \hline
        \textbf{Mô hình} & \textbf{Ưu điểm} & \textbf{Nhược điểm} \\
        \hline
        
        GPT-4o (OpenAI)
        & Lập luận phức tạp tốt, tool calling chính xác, hỗ trợ tiếng Việt tốt
        & Chi phí cao hơn GPT-4o mini \\
        \hline
        
        GPT-4o mini (OpenAI)
        & Chi phí thấp, tốc độ nhanh
        & Lập luận đa bước và tool calling kém chính xác hơn với các tình huống phức tạp \\
        \hline
        
        LLaMA 3.3 70B (Meta/Groq)
        & Tốc độ rất cao nhờ Groq LPU, chi phí thấp
        & Tool calling kém ổn định hơn GPT-4o trong một số tình huống \\
        \hline
    \end{tabular}
\end{table}

GPT-4o được lựa chọn vì khả năng lập luận đa bước và tool calling chính xác cao. Node Reason là thành phần điều phối trung tâm của agent, có thể thực hiện nhiều vòng gọi công cụ trong một lượt hội thoại, do đó chất lượng ra quyết định được ưu tiên hơn chi phí.

Để hỗ trợ tìm kiếm ngữ nghĩa tiếng Việt trong ChromaDB, hệ thống cần một mô hình embedding phù hợp với ngôn ngữ này. Các lựa chọn được xem xét:

\begin{table}[H]
    \centering
    \small
    \renewcommand{\arraystretch}{1.3}
    \begin{tabular}{|p{3cm}|p{5cm}|p{5cm}|}
        \hline
        \textbf{Mô hình} & \textbf{Ưu điểm} & \textbf{Nhược điểm} \\
        \hline
        
        text-embedding-3 (OpenAI)
        & Chất lượng cao, đa ngôn ngữ
        & Chi phí theo lượt gọi API, phụ thuộc dịch vụ bên ngoài \\
        \hline
        
        multilingual-e5-large (Microsoft)
        & Đa ngôn ngữ, chạy local
        & Chưa được tối ưu đặc biệt cho tiếng Việt \\
        \hline
        
        vietnamese-embedding (dangvantuan)
        & Được huấn luyện đặc biệt cho tiếng Việt, chạy local, miễn phí
        & Chỉ hỗ trợ tiếng Việt \\
        \hline
    \end{tabular}
\end{table}

Mô hình \texttt{dangvantuan/vietnamese-embedding} (\texttt{SentenceTransformer}) được lựa chọn vì được huấn luyện đặc biệt cho tiếng Việt, đảm bảo chất lượng biểu diễn ngữ nghĩa trong ngữ cảnh hội thoại tiếng Việt. Mô hình chạy hoàn toàn cục bộ, không phát sinh chi phí theo lượt truy vấn và không phụ thuộc vào dịch vụ bên ngoài.

\subsubsection{Công nghệ xây dựng API: FastAPI}

Hệ thống cần một framework để xây dựng API có hiệu năng cao và hỗ trợ tốt cho các tác vụ bất đồng bộ.

Các lựa chọn:

\begin{tabular}{|p{2.5cm}|p{5.5cm}|p{5.5cm}|}
\hline
Công nghệ & Ưu điểm & Nhược điểm \\
\hline
Flask & Đơn giản, dễ dùng & Không hỗ trợ async tốt \\
\hline
Django & Full-stack, nhiều tính năng & Nặng, dư thừa \\
\hline
FastAPI & Hiệu năng cao, hỗ trợ async, typing & Cộng đồng nhỏ hơn Django \\
\hline
\end{tabular}

FastAPI được lựa chọn do hỗ trợ bất đồng bộ, hiệu năng cao và tích hợp tốt với Python typing.
